% ============================================================
% Diagrama de Blocos — Controlador PID em Cascata
%
% Uso na figura:
%   \begin{figure}[H]
%     \centering
%     \caption{Arquitetura do controlador PID em cascata.}
%     \resizebox{\textwidth}{!}{%
%         % ============================================================
% Diagrama de Blocos — Controlador PID em Cascata
%
% Uso na figura:
%   \begin{figure}[H]
%     \centering
%     \caption{Arquitetura do controlador PID em cascata.}
%     \resizebox{\textwidth}{!}{%
%         % ============================================================
% Diagrama de Blocos — Controlador PID em Cascata
%
% Uso na figura:
%   \begin{figure}[H]
%     \centering
%     \caption{Arquitetura do controlador PID em cascata.}
%     \resizebox{\textwidth}{!}{%
%         % ============================================================
% Diagrama de Blocos — Controlador PID em Cascata
%
% Uso na figura:
%   \begin{figure}[H]
%     \centering
%     \caption{Arquitetura do controlador PID em cascata.}
%     \resizebox{\textwidth}{!}{%
%         \input{images/referencial_teorico/pid_cascata.tex}%
%     }
%     \label{fig:cascata}
%     \par\vspace{0.2cm}
%     \centering\footnotesize Fonte: Elaborado pelo autor (2025),
%     baseado em \citeonline{lopezsanchez2023}.
%   \end{figure}
% ============================================================

\begin{tikzpicture}[
  bloco/.style={
    rectangle, draw, thick,
    minimum width=2.6cm, minimum height=1.2cm,
    align=center, rounded corners=3pt, font=\small
  },
  blocomix/.style={
    rectangle, draw, thick,
    minimum width=1.5cm, minimum height=1.2cm,
    align=center, rounded corners=3pt, font=\small, fill=gray!12
  },
  soma/.style={
    circle, draw, thick,
    minimum size=0.6cm, inner sep=0pt, font=\normalsize
  },
  arr/.style={-{Stealth[length=5pt,width=4pt]}, thick},
  fb/.style={-{Stealth[length=5pt,width=4pt]}, thick, gray}
]

% ---- Junções de soma ----------------------------------------
\node[soma] at (0.0, 0) (s1) {$+$};
\node[soma] at (5.3, 0) (s2) {$+$};
\node[soma] at (10.6,0) (s3) {$+$};

% ---- Blocos -------------------------------------------------
\node[bloco, fill=violet!15] at (2.3,  0) (pp)
  {\textbf{PID Posi\c{c}\~ao}\\\footnotesize$\sim$50--100\,Hz};

\node[bloco, fill=teal!15]   at (7.6,  0) (pa)
  {\textbf{PID Atitude}\\\footnotesize$\sim$250--500\,Hz};

\node[bloco, fill=cyan!12]   at (12.9, 0) (pr)
  {\textbf{PID Taxa}\\\footnotesize$\sim$500--1000\,Hz};

\node[blocomix]              at (15.95,0) (mx) {\textbf{Mixing}};

\node[bloco, fill=orange!15] at (18.8, 0) (pl)
  {\textbf{Motores}\\\footnotesize e Planta};

% ---- Sinais negativos --------------------------------------
\node[font=\footnotesize] at ($(s1)+(0,-0.38)$) {$-$};
\node[font=\footnotesize] at ($(s2)+(0,-0.38)$) {$-$};
\node[font=\footnotesize] at ($(s3)+(0,-0.38)$) {$-$};

% ---- Entrada -----------------------------------------------
\draw[arr] (-1.4,0) -- (s1);
\node[font=\normalsize, above=0.85cm] at (-0.7, 0)
  {Ref.\ $(x_d,\,y_d,\,z_d)$};

% ---- Setas diretas ------------------------------------------
\draw[arr] (s1)      -- (pp);
\draw[arr] (pp.east) -- (s2.west);
\draw[arr] (s2)      -- (pa);
\draw[arr] (pa.east) -- (s3.west);
\draw[arr] (s3)      -- (pr);
\draw[arr] (pr.east) -- (mx.west);
\draw[arr] (mx.east) -- (pl.west);
\draw[arr] (pl.east) -- ++(0.6,0) coordinate (saida);

% ---- Rótulos intermediários — \normalsize acima dos blocos -
\node[font=\normalsize, above=0.85cm]
  at ($(pp.east)!0.5!(s2.west)$)
  {$\phi_d,\;\theta_d,\;T$};

\node[font=\normalsize, above=0.85cm]
  at ($(pa.east)!0.5!(s3.west)$)
  {$p_d,\;q_d,\;r_d$};

\node[font=\normalsize, above=0.85cm]
  at ($(pr.east)!0.5!(mx.west)$)
  {$\tau_\phi,\;\tau_\theta,\;\tau_\psi$};

% ---- Realimentações ----------------------------------------
\draw[fb] (saida) -- ++(0,-1.0) -| (s3.south);
\node[font=\small, black, anchor=west]
  at ($(s3.south)+(0.15,-1.02)$)
  {taxas $(p,\,q,\,r)$ --- girosc\'opio};

\draw[fb] (saida) -- ++(0,-1.9) -| (s2.south);
\node[font=\small, black, anchor=west]
  at ($(s2.south)+(0.15,-1.92)$)
  {\^Angulos $(\phi,\,\theta,\,\psi)$ --- IMU};

\draw[fb] (saida) -- ++(0,-2.8) -| (s1.south);
\node[font=\small, black, anchor=west]
  at ($(s1.south)+(0.15,-2.82)$)
  {Posi\c{c}\~ao $(x,\,y,\,z)$ --- fluxo \'optico / GPS};

\end{tikzpicture}%
%     }
%     \label{fig:cascata}
%     \par\vspace{0.2cm}
%     \centering\footnotesize Fonte: Elaborado pelo autor (2025),
%     baseado em \citeonline{lopezsanchez2023}.
%   \end{figure}
% ============================================================

\begin{tikzpicture}[
  bloco/.style={
    rectangle, draw, thick,
    minimum width=2.6cm, minimum height=1.2cm,
    align=center, rounded corners=3pt, font=\small
  },
  blocomix/.style={
    rectangle, draw, thick,
    minimum width=1.5cm, minimum height=1.2cm,
    align=center, rounded corners=3pt, font=\small, fill=gray!12
  },
  soma/.style={
    circle, draw, thick,
    minimum size=0.6cm, inner sep=0pt, font=\normalsize
  },
  arr/.style={-{Stealth[length=5pt,width=4pt]}, thick},
  fb/.style={-{Stealth[length=5pt,width=4pt]}, thick, gray}
]

% ---- Junções de soma ----------------------------------------
\node[soma] at (0.0, 0) (s1) {$+$};
\node[soma] at (5.3, 0) (s2) {$+$};
\node[soma] at (10.6,0) (s3) {$+$};

% ---- Blocos -------------------------------------------------
\node[bloco, fill=violet!15] at (2.3,  0) (pp)
  {\textbf{PID Posi\c{c}\~ao}\\\footnotesize$\sim$50--100\,Hz};

\node[bloco, fill=teal!15]   at (7.6,  0) (pa)
  {\textbf{PID Atitude}\\\footnotesize$\sim$250--500\,Hz};

\node[bloco, fill=cyan!12]   at (12.9, 0) (pr)
  {\textbf{PID Taxa}\\\footnotesize$\sim$500--1000\,Hz};

\node[blocomix]              at (15.95,0) (mx) {\textbf{Mixing}};

\node[bloco, fill=orange!15] at (18.8, 0) (pl)
  {\textbf{Motores}\\\footnotesize e Planta};

% ---- Sinais negativos --------------------------------------
\node[font=\footnotesize] at ($(s1)+(0,-0.38)$) {$-$};
\node[font=\footnotesize] at ($(s2)+(0,-0.38)$) {$-$};
\node[font=\footnotesize] at ($(s3)+(0,-0.38)$) {$-$};

% ---- Entrada -----------------------------------------------
\draw[arr] (-1.4,0) -- (s1);
\node[font=\normalsize, above=0.85cm] at (-0.7, 0)
  {Ref.\ $(x_d,\,y_d,\,z_d)$};

% ---- Setas diretas ------------------------------------------
\draw[arr] (s1)      -- (pp);
\draw[arr] (pp.east) -- (s2.west);
\draw[arr] (s2)      -- (pa);
\draw[arr] (pa.east) -- (s3.west);
\draw[arr] (s3)      -- (pr);
\draw[arr] (pr.east) -- (mx.west);
\draw[arr] (mx.east) -- (pl.west);
\draw[arr] (pl.east) -- ++(0.6,0) coordinate (saida);

% ---- Rótulos intermediários — \normalsize acima dos blocos -
\node[font=\normalsize, above=0.85cm]
  at ($(pp.east)!0.5!(s2.west)$)
  {$\phi_d,\;\theta_d,\;T$};

\node[font=\normalsize, above=0.85cm]
  at ($(pa.east)!0.5!(s3.west)$)
  {$p_d,\;q_d,\;r_d$};

\node[font=\normalsize, above=0.85cm]
  at ($(pr.east)!0.5!(mx.west)$)
  {$\tau_\phi,\;\tau_\theta,\;\tau_\psi$};

% ---- Realimentações ----------------------------------------
\draw[fb] (saida) -- ++(0,-1.0) -| (s3.south);
\node[font=\small, black, anchor=west]
  at ($(s3.south)+(0.15,-1.02)$)
  {taxas $(p,\,q,\,r)$ --- girosc\'opio};

\draw[fb] (saida) -- ++(0,-1.9) -| (s2.south);
\node[font=\small, black, anchor=west]
  at ($(s2.south)+(0.15,-1.92)$)
  {\^Angulos $(\phi,\,\theta,\,\psi)$ --- IMU};

\draw[fb] (saida) -- ++(0,-2.8) -| (s1.south);
\node[font=\small, black, anchor=west]
  at ($(s1.south)+(0.15,-2.82)$)
  {Posi\c{c}\~ao $(x,\,y,\,z)$ --- fluxo \'optico / GPS};

\end{tikzpicture}%
%     }
%     \label{fig:cascata}
%     \par\vspace{0.2cm}
%     \centering\footnotesize Fonte: Elaborado pelo autor (2025),
%     baseado em \citeonline{lopezsanchez2023}.
%   \end{figure}
% ============================================================

\begin{tikzpicture}[
  bloco/.style={
    rectangle, draw, thick,
    minimum width=2.6cm, minimum height=1.2cm,
    align=center, rounded corners=3pt, font=\small
  },
  blocomix/.style={
    rectangle, draw, thick,
    minimum width=1.5cm, minimum height=1.2cm,
    align=center, rounded corners=3pt, font=\small, fill=gray!12
  },
  soma/.style={
    circle, draw, thick,
    minimum size=0.6cm, inner sep=0pt, font=\normalsize
  },
  arr/.style={-{Stealth[length=5pt,width=4pt]}, thick},
  fb/.style={-{Stealth[length=5pt,width=4pt]}, thick, gray}
]

% ---- Junções de soma ----------------------------------------
\node[soma] at (0.0, 0) (s1) {$+$};
\node[soma] at (5.3, 0) (s2) {$+$};
\node[soma] at (10.6,0) (s3) {$+$};

% ---- Blocos -------------------------------------------------
\node[bloco, fill=violet!15] at (2.3,  0) (pp)
  {\textbf{PID Posi\c{c}\~ao}\\\footnotesize$\sim$50--100\,Hz};

\node[bloco, fill=teal!15]   at (7.6,  0) (pa)
  {\textbf{PID Atitude}\\\footnotesize$\sim$250--500\,Hz};

\node[bloco, fill=cyan!12]   at (12.9, 0) (pr)
  {\textbf{PID Taxa}\\\footnotesize$\sim$500--1000\,Hz};

\node[blocomix]              at (15.95,0) (mx) {\textbf{Mixing}};

\node[bloco, fill=orange!15] at (18.8, 0) (pl)
  {\textbf{Motores}\\\footnotesize e Planta};

% ---- Sinais negativos --------------------------------------
\node[font=\footnotesize] at ($(s1)+(0,-0.38)$) {$-$};
\node[font=\footnotesize] at ($(s2)+(0,-0.38)$) {$-$};
\node[font=\footnotesize] at ($(s3)+(0,-0.38)$) {$-$};

% ---- Entrada -----------------------------------------------
\draw[arr] (-1.4,0) -- (s1);
\node[font=\normalsize, above=0.85cm] at (-0.7, 0)
  {Ref.\ $(x_d,\,y_d,\,z_d)$};

% ---- Setas diretas ------------------------------------------
\draw[arr] (s1)      -- (pp);
\draw[arr] (pp.east) -- (s2.west);
\draw[arr] (s2)      -- (pa);
\draw[arr] (pa.east) -- (s3.west);
\draw[arr] (s3)      -- (pr);
\draw[arr] (pr.east) -- (mx.west);
\draw[arr] (mx.east) -- (pl.west);
\draw[arr] (pl.east) -- ++(0.6,0) coordinate (saida);

% ---- Rótulos intermediários — \normalsize acima dos blocos -
\node[font=\normalsize, above=0.85cm]
  at ($(pp.east)!0.5!(s2.west)$)
  {$\phi_d,\;\theta_d,\;T$};

\node[font=\normalsize, above=0.85cm]
  at ($(pa.east)!0.5!(s3.west)$)
  {$p_d,\;q_d,\;r_d$};

\node[font=\normalsize, above=0.85cm]
  at ($(pr.east)!0.5!(mx.west)$)
  {$\tau_\phi,\;\tau_\theta,\;\tau_\psi$};

% ---- Realimentações ----------------------------------------
\draw[fb] (saida) -- ++(0,-1.0) -| (s3.south);
\node[font=\small, black, anchor=west]
  at ($(s3.south)+(0.15,-1.02)$)
  {taxas $(p,\,q,\,r)$ --- girosc\'opio};

\draw[fb] (saida) -- ++(0,-1.9) -| (s2.south);
\node[font=\small, black, anchor=west]
  at ($(s2.south)+(0.15,-1.92)$)
  {\^Angulos $(\phi,\,\theta,\,\psi)$ --- IMU};

\draw[fb] (saida) -- ++(0,-2.8) -| (s1.south);
\node[font=\small, black, anchor=west]
  at ($(s1.south)+(0.15,-2.82)$)
  {Posi\c{c}\~ao $(x,\,y,\,z)$ --- fluxo \'optico / GPS};

\end{tikzpicture}%
%     }
%     \label{fig:cascata}
%     \par\vspace{0.2cm}
%     \centering\footnotesize Fonte: Elaborado pelo autor (2025),
%     baseado em \citeonline{lopezsanchez2023}.
%   \end{figure}
% ============================================================

\begin{tikzpicture}[
  bloco/.style={
    rectangle, draw, thick,
    minimum width=2.6cm, minimum height=1.2cm,
    align=center, rounded corners=3pt, font=\small
  },
  blocomix/.style={
    rectangle, draw, thick,
    minimum width=1.5cm, minimum height=1.2cm,
    align=center, rounded corners=3pt, font=\small, fill=gray!12
  },
  soma/.style={
    circle, draw, thick,
    minimum size=0.6cm, inner sep=0pt, font=\normalsize
  },
  arr/.style={-{Stealth[length=5pt,width=4pt]}, thick},
  fb/.style={-{Stealth[length=5pt,width=4pt]}, thick, gray}
]

% ---- Junções de soma ----------------------------------------
\node[soma] at (0.0, 0) (s1) {$+$};
\node[soma] at (5.3, 0) (s2) {$+$};
\node[soma] at (10.6,0) (s3) {$+$};

% ---- Blocos -------------------------------------------------
\node[bloco, fill=violet!15] at (2.3,  0) (pp)
  {\textbf{PID Posi\c{c}\~ao}\\\footnotesize$\sim$50--100\,Hz};

\node[bloco, fill=teal!15]   at (7.6,  0) (pa)
  {\textbf{PID Atitude}\\\footnotesize$\sim$250--500\,Hz};

\node[bloco, fill=cyan!12]   at (12.9, 0) (pr)
  {\textbf{PID Taxa}\\\footnotesize$\sim$500--1000\,Hz};

\node[blocomix]              at (15.95,0) (mx) {\textbf{Mixing}};

\node[bloco, fill=orange!15] at (18.8, 0) (pl)
  {\textbf{Motores}\\\footnotesize e Planta};

% ---- Sinais negativos --------------------------------------
\node[font=\footnotesize] at ($(s1)+(0,-0.38)$) {$-$};
\node[font=\footnotesize] at ($(s2)+(0,-0.38)$) {$-$};
\node[font=\footnotesize] at ($(s3)+(0,-0.38)$) {$-$};

% ---- Entrada -----------------------------------------------
\draw[arr] (-1.4,0) -- (s1);
\node[font=\normalsize, above=0.85cm] at (-0.7, 0)
  {Ref.\ $(x_d,\,y_d,\,z_d)$};

% ---- Setas diretas ------------------------------------------
\draw[arr] (s1)      -- (pp);
\draw[arr] (pp.east) -- (s2.west);
\draw[arr] (s2)      -- (pa);
\draw[arr] (pa.east) -- (s3.west);
\draw[arr] (s3)      -- (pr);
\draw[arr] (pr.east) -- (mx.west);
\draw[arr] (mx.east) -- (pl.west);
\draw[arr] (pl.east) -- ++(0.6,0) coordinate (saida);

% ---- Rótulos intermediários — \normalsize acima dos blocos -
\node[font=\normalsize, above=0.85cm]
  at ($(pp.east)!0.5!(s2.west)$)
  {$\phi_d,\;\theta_d,\;T$};

\node[font=\normalsize, above=0.85cm]
  at ($(pa.east)!0.5!(s3.west)$)
  {$p_d,\;q_d,\;r_d$};

\node[font=\normalsize, above=0.85cm]
  at ($(pr.east)!0.5!(mx.west)$)
  {$\tau_\phi,\;\tau_\theta,\;\tau_\psi$};

% ---- Realimentações ----------------------------------------
\draw[fb] (saida) -- ++(0,-1.0) -| (s3.south);
\node[font=\small, black, anchor=west]
  at ($(s3.south)+(0.15,-1.02)$)
  {taxas $(p,\,q,\,r)$ --- girosc\'opio};

\draw[fb] (saida) -- ++(0,-1.9) -| (s2.south);
\node[font=\small, black, anchor=west]
  at ($(s2.south)+(0.15,-1.92)$)
  {\^Angulos $(\phi,\,\theta,\,\psi)$ --- IMU};

\draw[fb] (saida) -- ++(0,-2.8) -| (s1.south);
\node[font=\small, black, anchor=west]
  at ($(s1.south)+(0.15,-2.82)$)
  {Posi\c{c}\~ao $(x,\,y,\,z)$ --- fluxo \'optico / GPS};

\end{tikzpicture}